% Generated draft: requires mathematical and source review.
\documentclass[11pt]{article}
\usepackage[margin=25mm]{geometry}
\usepackage{fontspec}
\setmainfont{Latin Modern Roman}
\usepackage{amsmath,amssymb,mathtools,unicode-math}
\setmathfont{Latin Modern Math}
\usepackage{enumitem,needspace,xurl,hyperref}
\setlength{\parindent}{0pt}
\setlength{\parskip}{5pt}
\setlength{\emergencystretch}{3em}
\newcommand{\N}{\mathbb{N}}
\newcommand{\Z}{\mathbb{Z}}
\newcommand{\Q}{\mathbb{Q}}
\newcommand{\R}{\mathbb{R}}
\newcommand{\C}{\mathbb{C}}
\newcommand{\F}{\mathbb{F}}
\newcommand{\A}{\mathbb{A}}
\newcommand{\PP}{\mathbb{P}}
\newcommand{\E}{\mathbb{E}}
\newcommand{\eps}{\varepsilon}
\newcommand{\dd}{\,\mathrm{d}}
\begin{document}
\section*{UnsolvedMath ID: 20001989}
\textbf{Source problem code:} catalogue entry
\subsection{Definitions and mathematical statement}
David Eisenbud's A.7 asks what Riemann--Roch, nodal plane curves, and residual intersection should mean for tropical plane curves. It suggests either line bundles/coherent sheaves or a geometric Brill--Noether approach. Classically, at a common point \(x\) of curves \(A,B,C\), with local ring \(\mathcal O_x\) on \(A\) and local equations \(f_B,f_C\), it records
\[
 ((f_B):f_C)=\{g\in\mathcal O_x:gf_C\in(f_B)\}
\]
as the residual ideal. Proper intersections and any containment needed for literal effective subtraction are not specified. At a smooth curve point multiplicities describe intersection divisors, but subtracting a larger multiplicity need not produce an effective residual cycle.

The source literally says ``a Gorenstein ring for any smooth curve, no matter how singular''. Its verification identifies this as internally contradictory and suggests replacing ``smooth'' by ``plane'', since a local plane-curve ring is a hypersurface quotient of a regular local ring. That replacement is an unverified inference, not adopted as an exact source correction. The record's clean-statement field contains this correction discussion rather than the actual full question.

The aid addresses the divisor question on a compact metric graph, not a tropical local colon-ideal construction. A divisor is a finite integral sum \(D=\sum n_p[p]\), effective when every \(n_p\ge0\), with \(\deg D=\sum n_p\). Linear equivalence is generated by principal divisors of continuous piecewise-linear functions with integer slopes. Its rank \(r(D)\) is \(-1\) if its class has no effective representative; otherwise it is the greatest \(r\ge0\) such that \(D-E\) is equivalent to an effective divisor for every effective \(E\) of degree \(r\). In the metric-graph Riemann--Roch framework,
\[
 r(D)-r(K-D)=\deg D+1-g,
 \qquad K=\sum_v(\operatorname{val}(v)-2)[v].
\]
This states the framework invoked by the aid, with valence-two subdivision immaterial.

For the compact circle skeleton \(\Gamma\) of a smooth tropical cubic, let its circumference be \(L>0\), choose an orientation and origin, and write \(s(p)\in\R/L\Z\) for arc-length coordinate. For a degree-zero divisor define
\[
 a(D)=\sum_p n_p s(p)\in\R/L\Z.
\]
The aid gives the exact genus-one rule
\[
 r(D)=
 \begin{cases}
 -1,&\deg D<0,\\
 \deg D-1,&\deg D>0,\\
 0,&\deg D=0\text{ and }a(D)=0,\\
 -1,&\deg D=0\text{ and }a(D)\ne0.
 \end{cases}
\]
For degree zero, \(a(D)=0\) is also the test for principality; it is independent of the chosen origin, and reversing orientation does not change its vanishing.

For stable-intersection divisors on the tropical cubic, restrict the following test to divisors \(D_B,D_C\) on \(\Gamma\), or explicitly push them to \(\Gamma\) by a specified skeleton retraction. Put \(D=D_B-D_C\). Three different notions must be separated:
\begin{enumerate}
\item A literal effective residual cycle requires coefficientwise \(D_B\ge D_C\).
\item An effective representative of the class exists exactly when \(r(D)\ge0\), as given by the displayed degree and coordinate test.
\item A principal difference requires \(\deg D=0\) and \(a(D)=0\).
\end{enumerate}
For degree zero an effective divisor itself must be zero, whereas a nonzero signed principal divisor can represent the effective zero class. The aid reports a length-13 example of this distinction, but does not give its points or coefficients. Retraction and divisor-class information do not construct a local tropical scheme or its colon ideal. The treatment of ordinary nodes, arbitrary singular tropical plane curves, and coherent-sheaf residuation therefore remains outside this reported genus-one partial result.
\subsection{Short English statement}
The question seeks tropical Riemann--Roch and a meaningful residual intersection for tropical plane curves. The aid gives an exact divisor-class test on a metric circle, distinguishing literal effectivity, effectivity up to linear equivalence, and principality. It does not construct a tropical analogue of the classical local colon ideal, and the proposed correction of a contradictory source sentence remains unverified.
\subsection{Sources}
\begin{itemize}
\item Local UnsolvedMath registry, permanent ID 20001989, exact source code catalogue entry; canonical \texttt{aim-geometry-notes.json}, \texttt{notes[326]}, tag \texttt{section}; AIM Workshop Problem Lists, Geometry, \emph{Amoebas and tropical geometry}, A.7, contributed by David Eisenbud. Created and updated 14 August 2026.
\item AIM source \url{https://aimath.org/WWN/amoebas/amoebas.pdf}. The supplied verification labels the statement \texttt{reconstructed\_unverified} and explicitly identifies the proposed plane-for-smooth replacement as inference. The PDF was not accessed.
\item The supplied aid invokes metric-graph Riemann--Roch but gives no author, theorem number, or full bibliographic citation for it. No authoritative reference or construction of tropical local colon ideals is supplied.
\item Attempt 1: \texttt{partial\_result}, problem status at run \texttt{partially\_solved}, rigor label \texttt{proved}, novelty confidence \texttt{low}; classification \texttt{PARTIAL-PROGRESS}, registry status \texttt{partially\_solved}. Candidate: genus-one residual test and a reported length-13 example. These are dated source labels and claims requiring review.
\item Ulam AI / UnsolvedMath Contributors; CC BY 4.0. Machine-generated research aid, not peer reviewed; independently verify before citation.
\end{itemize}
\end{document}
